%% guide-fosse-mariannes — colonne d'eau, du niveau de la mer au fond de la fosse.
%% Dégradé solideB!15 (surface) -> solideB!55 (fond) ; plancher océanique hachuré en V.
%% Le bathyscaphe est posé au fond ; les trois flèches rouges rappellent que la pression
%% s'exerce dans toutes les directions et croît avec la profondeur.
%% RÈGLE : aucune pression calculée (ni Delta P, ni P abs) — seules figurent les données
%% de l'énoncé : P atm = 1 013 hPa, rho = 1 030 kg/m3, h = 10 994 m, mesure de 1960.
\documentclass[tikz,border=4pt]{standalone}
\usepackage{liaisons}
\usetikzlibrary{decorations.pathmorphing}
\begin{document}
\begin{tikzpicture}[x=1cm,y=1cm,
  cote/.style={{Stealth[length=4.5pt]}-{Stealth[length=4.5pt]}, line width=0.9pt, draw=solideD},
  attache/.style={line width=0.4pt, draw=black!45},
  pousse/.style={-{Stealth[length=5pt]}, line width=1.3pt, draw=solideA},
  lab/.style={font=\small, inner sep=1.5pt},
  pet/.style={font=\scriptsize, inner sep=1.5pt}]

\useasboundingbox (0.30,-1.75) rectangle (7.90,4.15);
\def\xg{1.70}\def\xd{6.60}\def\ys{3.30}   % bords et surface libre de la colonne d'eau

%% ---- profil du plancher océanique (fosse en V) --------------------------------
\def\fond{(\xd,1.60) -- (5.70,1.10) -- (5.00,0.50) -- (4.35,0.14) -- (3.95,0.06)
          -- (3.45,0.32) -- (2.90,0.92) -- (2.25,1.50) -- (\xg,1.86)}

%% ---- l'eau : dégradé du haut vers le bas ---------------------------------------
\shade[top color=solideB!15, bottom color=solideB!60]
  (\xg,\ys) -- (\xd,\ys) -- \fond -- cycle;
\draw[line width=1.1pt, solideB, decorate,
      decoration={snake, amplitude=0.7pt, segment length=7pt}] (\xg,\ys) -- (\xd,\ys);
\draw[line width=0.8pt, black!55] (\xg,\ys) -- (\xg,0.0) (\xd,\ys) -- (\xd,0.0);

%% ---- fond hachuré ----------------------------------------------------------------
\hachures{\fond -- (\xg,-0.20) -- (\xd,-0.20) -- cycle}

%% ---- l'air et la pression atmosphérique --------------------------------------------
\node[lab, anchor=south] at (4.15,3.62) {AIR --- $P_{\text{atm}} = 1\,013$ hPa};
\foreach \x in {2.55,4.15,5.75} {\draw[-{Stealth[length=4pt]}, line width=0.8pt, solideD]
  (\x,3.55) -- (\x,3.40);}
\node[pet, text=solideB!80!black, anchor=north west] at (\xg+0.12,\ys-0.10)
  {eau de mer, $\rho = 1\,030$ kg$\cdot$m$^{-3}$};

%% ---- cote de la profondeur ------------------------------------------------------------
\draw[attache] (1.00,\ys) -- (\xg-0.04,\ys);
\draw[attache] (1.00,0.06) -- (3.90,0.06);
\draw[cote] (1.12,\ys) -- (1.12,0.06);
\node[lab, text=solideD, rotate=90] at (0.90,1.70) {$h = 10\,994$ m};

%% ---- bathyscaphe posé au fond -----------------------------------------------------------
\begin{scope}[shift={(3.95,0.98)}]
  \draw[line width=1pt, fill=black!25] (-0.62,-0.02) .. controls (-0.62,0.34) and (0.62,0.34)
       .. (0.62,-0.02) .. controls (0.62,-0.24) and (-0.62,-0.24) .. (-0.62,-0.02) -- cycle;
  \draw[line width=1pt, fill=black!45] (0,-0.30) circle (0.24);
  \draw[line width=0.7pt, fill=white] (0.12,-0.30) circle (0.06);
  \draw[line width=0.9pt] (-0.30,0.28) -- (-0.30,0.46) (0.30,0.28) -- (0.30,0.46);
\end{scope}
\node[pet, anchor=south, align=center, fill=white, draw=black!45, rounded corners=2pt,
      inner sep=2.5pt] at (3.05,1.98) {mesure de 1960 :\\\textbf{1 087 bars}};
\draw[line width=0.4pt, black!45] (3.35,1.95) -- (3.62,1.40);

%% ---- la pression s'exerce dans toutes les directions -------------------------------------
\draw[pousse] (3.95,1.86) -- (3.95,1.52);
\draw[pousse] (2.50,1.00) -- (3.18,1.00);
\draw[pousse] (5.40,0.86) -- (4.62,0.86);
\node[pet, text=solideA, anchor=north west, align=left] at (4.90,2.80)
  {la pression\\s'exerce dans\\toutes\\les directions};

%% ---- rappel encadré des relations ---------------------------------------------------------
\node[anchor=north, align=center, font=\scriptsize, draw=black!55, rounded corners=2pt,
      inner sep=4pt, fill=white] at (4.10,-0.45)
  {$\Delta P = \rho \, g \, h$ \qquad $P_{\text{abs}} = P_{\text{atm}} + \Delta P$\\[1pt]
   $P$ en Pa \; ; \; $\rho$ en kg$\cdot$m$^{-3}$ \; ; \; $g$ en N$\cdot$kg$^{-1}$ \; ; \; $h$ en m};
\end{tikzpicture}
\end{document}
